\chapter{Ultrasound}
\section*{What Is This Test?} Ultrasound uses high-frequency sound waves to create images of organs, tissues, blood flow, and developing pregnancies.\section*{Alternative Names} Sonography; diagnostic ultrasound; ultrasonography.\section*{Principle} Echoes from tissue boundaries are converted into images; Doppler evaluates motion and blood flow.\section*{Why This Test Is Ordered} It evaluates abdominal, pelvic, thyroid, breast, vascular, musculoskeletal, cardiac, and fetal conditions.\section*{Primary vs Secondary Test} It is a common noninvasive primary imaging study; CT, MRI, or biopsy may be secondary.\section*{How This Test Is Performed} Gel and a transducer are moved over the skin; some studies use an internal probe.\section*{What Preparation Is Needed} Fasting or a full bladder may be requested depending on the body region.\section*{Understanding Results} Findings depend on anatomy, operator technique, body habitus, and the clinical question; a normal study does not exclude all disease.\section*{Follow-up Tests} Repeat ultrasound, CT, MRI, laboratory testing, or biopsy may follow.\section*{References} \begin{enumerate}\item MedlinePlus. Ultrasound.\item American Institute of Ultrasound in Medicine. Ultrasound practice guidance.\end{enumerate}\clearpage\section*{Understanding Results and Follow-up}
The laboratory report should identify the specimen, method, units, reference interval or decision limit, and whether the result is positive, negative, abnormal, or inconclusive. Reference intervals vary with age, sex, pregnancy, specimen type, assay, and laboratory; a result outside a range is not automatically a diagnosis, and a result within a range does not always exclude disease. Discuss unexpected results with the ordering clinician, who can decide whether repeat or confirmatory testing, treatment, or referral is appropriate. Do not change medicines or delay urgent care based on an isolated result.
\section*{Regulatory Status and Authorization Date}This chapter describes a test category, not a specific commercial product. FDA, CDSCO, EU IVDR, and MHRA approval or registration dates therefore cannot be assigned without the assay manufacturer, product name, instrument, and intended use. For laboratory-developed tests, an individual agency approval date may not exist. See the Preface for the interpretation of regulatory status.
\clearpage
